%! Setting Up a Test for a Population Mean — Unit 7, Lesson 7.4 — Exit Ticket (Answer Key)
\documentclass[10pt]{article}
\usepackage{apstats-article}
\usepackage{apstats-key}   % pulls in -boxes; do NOT also load -boxes
\newcommand{\TallMath}[1]{$\displaystyle #1\rule[-1.4em]{0pt}{3.2em}$}

\begin{document}

\pageheader{Unit 7, Lesson 7.4}{Exit Ticket --- Answer Key}
\namedateperiod

\vspace{0.15in}

A physician believes the mean systolic blood pressure of adults in her practice is
\emph{higher} than the national average of 120 mmHg. She randomly selects 18 patients
and records their systolic blood pressure. A dotplot of the data shows a roughly symmetric
distribution with one mild outlier.

\begin{enumerate}
  \item Name the appropriate inference method and give the degrees of freedom.

        Method: \ans{one-sample $t$-test for $\mu$} \quad $df = $ \ans{17}

  \item State hypotheses.

        $H_0$: \ans{$\mu = 120$} \quad $H_a$: \ans{$\mu > 120$}

  \item Verify conditions. \textit{Be specific --- use numbers and describe what you see in the data.}
        \begin{itemize}
          \item Random: \ansline{The 18 patients were randomly selected --- Random condition met.}
          \item 10\% rule: \ansline{$18 \le 10\%$ of all patients in the practice (assuming more than 180 patients) --- met.}
          \item Normal ($n < 30$): \ansline{$n = 18 < 30$; dotplot shows roughly symmetric with one mild outlier. A single mild outlier in a symmetric distribution is generally acceptable --- Normal condition is borderline but likely met; proceed with caution.}

                Can the test proceed? \ans{Yes, with caution} \quad Explain: \ansline{The distribution is roughly symmetric; one mild outlier is not a severe violation. Proceed, but note the limitation.}
        \end{itemize}
\end{enumerate}

\begin{teachernote}
Use the ``mild outlier'' scenario to spark discussion: one clear outlier in a sample of 18
may still be acceptable if the rest of the distribution is symmetric. A severe outlier or
strong skewness would be cause to not proceed. Encourage students to describe both what they
see and their judgment about whether the condition is met.
\end{teachernote}

\end{document}
